\section{Related work}
\label{sec:related}

Probe sits at the intersection of four literatures: LLM-based peer review, AI-assisted research-design tooling, multi-agent systems for scientific work, and the autonomy-by-design tradition in critical HCI. Each cluster motivates specific architectural choices and clarifies what Probe is (a triage tool that defers to humans) versus what it refuses to be (an autonomous scientist).

\subsection{LLM-based peer review and reviewer simulation}

Work on LLMs as reviewers runs from complementary feedback \citep{liang2024large} through structured simulation \citep{jin2024agentreview, darcy2024marg, zhu2025deepreview, bougie2025gar} to full reviewer replacement in AI-Scientist-style pipelines. Probe borrows the role-specialized reviewer panel and the disagreement-preserving meta-review from this line but refuses the replacement framing; it inherits Liang's ``complement, not replace'' conclusion as the structural invariant that every guidebook must carry a \texttt{[HUMAN\_REQUIRED]} tag. AgentReview supplies the empirical warrant for role-specialized personas: introducing a single biased reviewer flips the accept/reject decision on 37.1\% of papers \citep{jin2024agentreview}. Probe treats this finding not as a bug to average away but as a signal to preserve through the meta-review rather than collapsing to a single consensus score. MARG \citep{darcy2024marg} decomposes review into aspect-specific sub-agents and reduces the rate of generic comments from 60\% (GPT-4 baseline) to 29\%; this is the architectural antecedent for Probe's Stage 7 split across methodologist, accessibility, and novelty reviewers. DeepReview \citep{zhu2025deepreview} shows that structured, literature-grounded reasoning chains reduce hallucinated critique, which is the motivation for Probe's forbidden-phrase linter and its provenance-tagged criticisms. GAR \citep{bougie2025gar} and ResearchAgent \citep{baek2025researchagent} demonstrate that reviewer agents can produce human-comparable feedback at scale, but the Dagstuhl-derived survey of NLP for peer review \citep{kuznetsov2024nlppeerreview} catalogues why full automation remains premature: data licensing, operationalization, and ethics problems that Probe addresses by scoping to early-stage triage rather than submission review. The citation-fabrication literature establishes the concrete failure mode that Probe's provenance linter and \texttt{corpus/source\_cards/} allowed-claims whitelist exist to prevent.

\subsection{Research-design and study-planning tools with AI assistance}

The immediate neighbors are systems that use LLMs to help researchers develop research questions \citep{liu2024coquest}, research ideas \citep{si2024novelideas, baek2025researchagent}, or full papers \citep{lu2024aiscientist, schmidgall2025agentlab}. None generates adversarially-reviewed, provenance-labeled study guidebooks for human workshops --- the gap Probe fills. Si, Yang, and Hashimoto's blind review with 100+ NLP researchers \citep{si2024novelideas} established the first statistically-significant finding that LLM-generated ideas are rated more novel but less feasible than expert ideas, which Probe treats as the motivation for the Stage-6 audit: if machine-generated study designs over-sell novelty and under-check feasibility, a dedicated feasibility audit with binding verdicts becomes load-bearing. ResearchAgent \citep{baek2025researchagent} introduces the core loop of LLM ideation with a reviewing-agent feedback stage, which Probe's eight-stage pipeline extends by adding explicit provenance, three divergent branches instead of one iterated idea, and a forbidden-phrase linter that rejects evidence language in Probe's own voice. CoQuest \citep{liu2024coquest} focuses on breadth-first versus depth-first ideation UX and finds that AI processing delays increase researcher reflection, which Probe operationalizes by running three branches to completion in parallel rather than returning a single best answer. On the cultural-probes side, Gaver et al.\ \citep{gaver1999probes} and Sanders and Stappers \citep{sanders2014codesign} supply the ``rehearsal, not evidence'' framing that Probe inherits verbatim; the \texttt{[SIMULATION\_REHEARSAL]} provenance tag exists because these authors explicitly rejected the generalizable-data reading of probe work. Jakesch et al.'s demonstration that opinionated language models shift user views \citep{jakesch2023cowriting} is the empirical grounding for Probe's Legibility axis: an AI that quietly reshapes the researcher's framing is the specific capture risk the self-audit is designed to flag.

\subsection{Multi-agent systems for complex scientific and creative tasks}

Probe's eight-stage pipeline with parallel branches sits inside a dense landscape of multi-agent frameworks, which divide along two axes: academic papers defining patterns and open-source pipelines that commit to a specific end-to-end task. CAMEL \citep{li2023camel} established inception-prompted role-play as a core pattern; AutoGen \citep{wu2024autogen} generalized this to composable conversation patterns; MetaGPT \citep{hong2024metagpt} encoded standardized operating procedures as the key to reducing inter-agent error --- Probe's per-stage typed schema validation and single-repair-pass policy derive from this SOP discipline. Multi-agent debate \citep{du2024debate} shows parallel agents with different viewpoints reduce hallucination; Probe's three-branch divergent ideation applies this pattern at the research-design level rather than the answer-refinement level.

Among open-source systems, AI-Scientist \citep{lu2024aiscientist} and Agent Laboratory \citep{schmidgall2025agentlab} are the closest full-pipeline analogues. An independent evaluation by Beel et al.\ \citep{beel2025aiscientistevaluation} found AI-Scientist's literature review yields ``poor novelty assessments'' and that 5 of 12 proposed experiments (42\%) failed on coding errors. MLR-Bench \citep{chen2025mlrbench} evaluates AI-Scientist v2 alongside its own MLR-Agent and reports that coding agents produce fabricated or invalidated experimental results in roughly 80\% of tasks --- a failure mode both systems share. Probe's response is structural: it does not run experiments or claim novelty without human verification; its Managed Agents deep audit uses bash and grep to measure quantitative claims against the repository itself, producing \texttt{[TOOL\_VERIFIED]} tags only for claims the audit could mechanically check. Karpathy's \texttt{autoresearch} \citep{karpathy2025autoresearch} is the stylistic antecedent for Probe's \texttt{PROBE.md} single-file orchestration contract; NousResearch's \texttt{autonovel} \citep{nousresearch2025autonovel} is the antecedent for the multi-stage critique-and-cut pipeline and the explicit anti-slop forbidden-phrase discipline; \texttt{feynman} \citep{companion2025feynman} is the closest design sibling on provenance, with every claim link-grounded to source URLs. Probe differs from all three by targeting HCI study design rather than novels, ML experiments, or literature reviews, and by enforcing capture-risk audits as a first-class stage rather than a post-hoc evaluation.

\subsection{Autonomy-by-design and capture-risk frameworks}

Probe's four-axis capture-risk audit (Capacity, Agency, Exit, Legibility) and its 16 named patterns inherit from a theoretical stack that runs from Illich's conviviality \citep{illich1973conviviality} through Kafer's political-relational disability model \citep{kafer2013feminist} to Buijsman, Carter, and Bermúdez's domain-specific-autonomy framework \citep{buijsman2025autonomy}. Illich's argument that tools past a scale threshold reverse from serving human purposes to conditioning humans to serve the tool is the Exit axis: Probe asks, for every study design it triages, whether the intervention reduces the researcher's or participant's space of available action once use stops. Kafer's critique of the curative imaginary is encoded in the \texttt{agency.paternalistic\_default} pattern and in the accessibility-advocate reviewer persona, which flags extractive framings even when the agents themselves produce them. Buijsman et al.'s analysis of AI decision-support as eroding ``skilled competence and authentic value-formation'' via absent failure indicators is the direct warrant for the audit's structure: each of the 16 patterns specifies a quoted-evidence template and a defeater mechanism, matching the design moves Buijsman et al.\ propose \citep{buijsman2025autonomy}.

Adjacent work in 2023--2025 extends this stack. Costanza-Chock's \emph{Design Justice} \citep{costanzachock2020designjustice} supplies the structural framing --- ``Nothing About Us Without Us'' as design constraint --- that motivates Probe's insistence on the \texttt{[HUMAN\_REQUIRED]} tag and on running a dedicated accessibility-advocate reviewer persona. The open question Probe positions itself against is stated most sharply in Buijsman et al.: how do we get the productivity benefits of AI decision-support without the erosion of domain-specific autonomy? Probe's answer is not to claim the erosion can be eliminated but to make it legible --- to produce, for every run, a provenance trail a researcher can inspect to see where the machine's framing shaped the guidebook.

\subsection{Direct-overlap check}

Among 2024--2025 work, the closest architectural sibling is ResearchAgent \citep{baek2025researchagent}; the closest venue-and-domain sibling is CoQuest \citep{liu2024coquest}; the closest fuller-pipeline analogues are AI-Scientist and Agent Laboratory \citep{lu2024aiscientist, schmidgall2025agentlab}. None of them (a) generates divergent HCI study guidebooks, (b) runs multi-persona adversarial review with blocking verdicts, (c) enforces provenance-labeled claims with a forbidden-phrase linter, and (d) applies an autonomy-by-design capture-risk audit. Probe appears to be the first system to combine these four elements. We flag this as a claim about scope and architectural combination rather than about individual-component novelty: each element has precedent in the literature; the combination does not.
